% =====================================================================
%  Analisis de contexto del replanteo
%
%  Reune lo que se sabe del problema antes de redisenar: lo que pide el
%  pliego, como se toman los examenes en la universidad, de donde salen
%  los datos, que se decidio y que falta decidir. Es el respaldo de las
%  decisiones de ../README.md.
%
%  Autocontenido: solo depende del logo de ../../guia/fichas/logo.
%  Compilar con ./compilar.sh desde esta carpeta.
% =====================================================================
\documentclass[11pt,letterpaper]{article}
\usepackage[utf8]{inputenc}
\usepackage[T1]{fontenc}
\usepackage[spanish,es-noquoting]{babel}
\usepackage[margin=2.2cm,headheight=24pt,headsep=10pt]{geometry}
\usepackage{array}
\usepackage{longtable}
\usepackage{graphicx}
\usepackage{fancyhdr}
\usepackage{parskip}
\usepackage{enumitem}
\usepackage{xcolor}
\usepackage[hidelinks]{hyperref}

\newcolumntype{L}[1]{>{\raggedright\arraybackslash}p{#1}}
\newcolumntype{C}[1]{>{\centering\arraybackslash}p{#1}}
\setlength{\emergencystretch}{3em}
\renewcommand{\arraystretch}{1.25}
\setlist{itemsep=2pt,topsep=3pt}

\newcommand{\Logo}{%
  \IfFileExists{../../guia/fichas/logo/nexasoft_horizontal.png}{%
    \includegraphics[width=2.5cm]{../../guia/fichas/logo/nexasoft_horizontal}}{\textbf{NexaSoft}}}

\pagestyle{fancy}
\fancyhf{}
\fancyhead[L]{\raisebox{-3pt}{\Logo}}
\fancyhead[R]{\footnotesize Análisis de contexto del replanteo}
\fancyfoot[C]{\footnotesize \thepage}
\renewcommand{\headrulewidth}{0.4pt}

% Origen de cada afirmacion: lo verificado contra una fuente, lo relatado
% por el equipo y lo que es propuesta todavia sin confirmar.
\newcommand{\relatado}{\textcolor{gray}{\footnotesize\,[relato del equipo]}}
\newcommand{\propuesto}{\textcolor{gray}{\footnotesize\,[propuesto]}}

\begin{document}

\begin{center}
{\Large\bfseries Análisis de contexto del replanteo}\\[4pt]
{\large Sistema de control de ingreso a exámenes masivos}\\[6pt]
{\small NexaSoft S.R.L. \,$\cdot$\, Proyecto guía \,$\cdot$\, 4 de octubre de 2026}
\end{center}

\section{Propósito}

Este documento reúne lo que se sabe del problema antes de rediseñar el sistema: qué pide
el pliego, cómo se toman en la práctica los exámenes en la Universidad Mayor de San Simón,
de dónde pueden salir los datos y qué se decidió a partir de eso. Es el respaldo de las
decisiones del replanteo; el diseño resultante está en \texttt{docs/replanteamiento}.

Cada afirmación lleva su origen. Lo que no tiene marca se verificó contra una fuente: el
pliego, la transcripción de una revisión, los datos de GENDA o una página pública de la
universidad. La marca \relatado\ señala lo que describió el equipo desde su experiencia
como estudiantes, sin documento que lo respalde. La marca \propuesto\ señala lo que
todavía no se confirmó.

\section{De dónde sale el replanteo}

En la octava revisión, del 29 de septiembre de 2026, la consultora probó el sistema como
administradora y se puso después en el lugar de una docente: con un examen de su materia
el 8 de octubre, quería la lista de sus estudiantes para generarles el código QR. El
sistema no podía dársela. Sus observaciones:

\begin{itemize}
  \item No se sabe de qué materia, grupo ni docente es cada estudiante. El padrón es una
        lista única y el administrador sube estudiantes sin separarlos.
  \item Hay que preguntarse para qué sirve cada dato que se guarda y qué consultas se
        van a necesitar después.
  \item La autoinscripción con un código de curso publicado deja entrar intrusos. Lo
        vivió en la virtualidad: un grupo de 180 terminó con más de 300 inscritos.
  \item El código QR tiene que ser por estudiante y no predecible.
  \item Usar la oferta académica real, obtenida de las páginas de las facultades, le
        pareció valioso por ser genérico. Pidió realismo con el tiempo: con dos o tres
        facultades alcanza para probar.
  \item La interfaz debe mostrar quién inició sesión y con qué rol, marcar los campos
        obligatorios, validar mientras se escribe y no hacer teclear el dominio del correo.
  \item Dudó del rol de responsable académico.
\end{itemize}

\section{Lo que exige el pliego}

El pliego describe un sistema para la «gestión, validación y control del ingreso de
estudiantes durante la realización de exámenes masivos» en una institución de educación
superior. De sus dieciséis requerimientos, estos condicionan el replanteo:

\begin{center}
\begin{longtable}{@{}C{1.1cm}L{8.2cm}L{6.4cm}@{}}
\hline
\textbf{Req.} & \textbf{Qué pide} & \textbf{Qué implica} \\
\hline
\endhead
1 & Registrar los datos de los estudiantes que participarán en los exámenes & Hace falta un padrón \\
2 & Registrar exámenes con asignatura, fecha, hora, duración y ambiente & El examen es de una asignatura \\
3 & Asociar los habilitados a un examen y saber por qué alguien no lo está & Habilitación con motivo \\
4 & Verificar la identidad por código universitario, documento o código QR & El estudiante no opera el sistema: se lo identifica \\
6 y 7 & Comprobar la habilitación y registrar la hora de ingreso & La verificación ocurre en la puerta \\
9 & Registrar el intento de ingreso de un estudiante «no registrado en la materia» & El sistema tiene que saber quién está inscrito en cada materia \\
10 & Verificar el ambiente en el que ingresó & El estudiante tiene un aula asignada \\
16 & Administrar tipos de usuario: «por ejemplo, administradores, docentes y personal encargado» & Los estudiantes no tienen cuenta. El responsable académico no figura \\
\hline
\end{longtable}
\end{center}

Conclusiones sobre el alcance:

\begin{itemize}
  \item Se trata de los exámenes de las materias durante la gestión: parciales, finales y
        segunda instancia. No del examen de ingreso a la universidad, donde no hay
        materia, grupo ni código universitario.
  \item «Masivo» se refiere a las materias grandes, en las que un examen junta varios
        grupos en varias aulas a la vez.
  \item Un enfoque opcional, al estilo de un aula virtual que cada docente usa si quiere,
        no cumple el requerimiento 9: sin lista confiable de inscritos no se puede saber
        quién no corresponde.
\end{itemize}

\section{Cómo se toman los exámenes en la universidad}

\subsection{Quién atiende el examen}

No existe un cargo de «personal de control». Quien controla el ingreso es el docente o
los auxiliares que él lleva. Se dan cuatro situaciones: \relatado

\begin{center}
\begin{tabular}{@{}L{4.6cm}L{11.2cm}@{}}
\hline
\textbf{Situación} & \textbf{Qué ocurre} \\
\hline
El docente solo & Toma el examen y hace todo el control por su cuenta \\
El docente con auxiliares & Lo ayudan estudiantes que trae para la ocasión, o el auxiliar que ya tiene asignado a su grupo. Controlan junto con él \\
Un auxiliar a cargo & El docente no puede asistir y envía a un auxiliar a tomar el examen. Es excepcional \\
Varias aulas & Los estudiantes son demasiados para un aula. El docente ocupa más de una, asigna un auxiliar a cada una y va pasando por todas \\
\hline
\end{tabular}
\end{center}

\subsection{Exámenes que comparten día, hora y aula}

\begin{itemize}
  \item \textbf{Docentes de una misma materia.} Los docentes de distintos grupos de una
        asignatura toman el examen a la vez y pueden compartir un aula. \relatado
  \item \textbf{Un docente con varias materias.} Un docente que dicta, por ejemplo,
        Álgebra I, Estadística y Cálculo puede citar a todos sus grupos el mismo día en la
        misma aula y tomarles el examen de forma simultánea. \relatado
  \item \textbf{El auditorio.} En la Facultad de Ciencias y Tecnología el auditorio es el
        aula más grande y es donde suelen rendir varios grupos juntos. \relatado\ En los
        horarios de la facultad figuran tres: \texttt{AUDCIV}, \texttt{AUDIND} y
        \texttt{AUDINF}. No están ubicados en el mapa de GENDA. Humanidades tiene los
        suyos ubicados: \texttt{AUD.1}, \texttt{AUD.2} y \texttt{AUD.P}.
\end{itemize}

\subsection{El estudiante}

\begin{itemize}
  \item Cursa unas seis materias por semestre, o por año en las carreras anuales.
        \relatado\ Si cada docente cargara a sus estudiantes por separado, el mismo
        estudiante se repetiría seis veces.
  \item No hay una forma segura de que se dé de alta por su cuenta: es la puerta a
        perfiles falsos y a inscripciones en materias ajenas.
  \item Los estudiantes de primer semestre no saben dónde quedan las aulas. Lo comentó la
        consultora a propósito del mapa: preguntan en la puerta dónde está el auditorio.
\end{itemize}

\section{Organización académica}

\subsection{Períodos}

La universidad numera sus períodos. Lo que se pudo establecer con el cronograma de la
WebSIS y el calendario de la Facultad de Ciencias y Tecnología:

\begin{center}
\begin{tabular}{@{}C{1.6cm}L{4.4cm}L{9.8cm}@{}}
\hline
\textbf{Código} & \textbf{Período} & \textbf{Observación} \\
\hline
1 & Semestre 1 & En 2026 terminó el 4 de julio (FCyT) \\
2 & Semestre 2 & 72 carreras lo tienen vigente \\
0 & Anual & 14 carreras: Arquitectura, Medicina, Odontología, Derecho y otras \\
4 & Curso de invierno & FCyT, del 9 al 31 de julio de 2026 \\
\hline
\end{tabular}
\end{center}

Al 4 de octubre de 2026 están vigentes a la vez \texttt{2/2026} y \texttt{0/2026}. El
código del curso de verano no se confirmó.

Fechas de la FCyT para \texttt{2/2026}: inicio el 10 de agosto; primeros parciales del 12
al 31 de octubre; segundos parciales del 7 al 17 de diciembre; finales y segunda
instancia del 19 al 24 de diciembre; cierre el 26 de diciembre. También hay exámenes de
mesa, del 10 al 17 de agosto.

\subsection{Facultades, materias y docentes}

\begin{itemize}
  \item Una asignatura tiene varios grupos, cada uno con su docente, su horario y sus
        aulas. Introducción a la Programación tiene ocho grupos en \texttt{2/2026};
        Física General, treinta.
  \item Un docente puede dictar en más de una facultad. En la oferta \texttt{2/2026} de
        las cuatro facultades analizadas hay 849 docentes y 29 están en esa situación;
        uno dicta en tres.
  \item Hay grupos sin docente designado al publicarse el horario.
  \item Los grupos pueden tener auxiliar titular. Figura en el horario publicado, aunque
        GENDA hoy lo descarta al procesarlo.
\end{itemize}

\section{De dónde pueden salir los datos}

\subsection{GENDA}

GENDA es un proyecto anterior del equipo que recopila la oferta académica publicada por
las facultades. Lo que tiene:

\begin{center}
\begin{tabular}{@{}L{4.2cm}L{11.6cm}@{}}
\hline
\textbf{Dato} & \textbf{Qué hay} \\
\hline
Horarios & Carrera, nivel, asignatura con código, grupo, docente, días, horas y aula. Cuatro facultades con datos de \texttt{2/2026}: Ciencias y Tecnología, Economía, Humanidades y Arquitectura \\
Pensum & 18 facultades, 105 carreras y 4.574 materias, con el régimen de cada carrera \\
Aulas & 237 aulas ubicadas en 56 edificios de esas cuatro facultades, con piso y coordenadas \\
Docentes & Solo el nombre, tomado de cada grupo. Sin documento ni correo \\
Colores & Un color por facultad, ya definido \\
\hline
\end{tabular}
\end{center}

Lo que no tiene: estudiantes, inscripciones, capacidad de las aulas y correo de los
docentes. Los datos de Derecho y Odontología son de \texttt{1/2025}, y Derecho no tiene
fuente web.

\subsection{Fuentes públicas de la universidad}

\begin{itemize}
  \item \textbf{Cronograma de inscripciones de la WebSIS.} Una página pública con el
        período vigente de cada una de las 92 carreras, con su código. Los códigos de
        carrera coinciden con los de GENDA. La fecha que publica es la de apertura de
        inscripciones, no la de inicio de clases, y trae planes cerrados con períodos
        viejos.
  \item \textbf{Calendario académico de la FCyT.} Un PDF con texto, una fila por día, con
        el inicio, el fin y las ventanas de exámenes.
  \item \textbf{Calendario de Economía.} Publicado como PDF e imagen; no se comprobó si
        se puede leer. De las demás facultades no se encontró calendario.
\end{itemize}

Los enlaces y la forma de leer cada fuente están en \texttt{database/umss/fuentes.json}.

\subsection{Listas de inscritos}

No son públicas. El docente descarga la lista de su grupo desde su área de la WebSIS.
\relatado\ No se sabe a qué información tiene acceso la administración: se asume que
puede obtener las inscripciones de una facultad, sin haberlo confirmado. Falta conseguir
un archivo de ejemplo para conocer su formato.

\section{Actores}

\begin{center}
\begin{tabular}{@{}L{3.4cm}L{12.4cm}@{}}
\hline
\textbf{Actor} & \textbf{Papel} \\
\hline
Administrador & Prepara cada período: importa la oferta, activa las cuentas de los docentes y verifica el padrón. Ve toda la universidad, por facultad. Se encarga de los datos del sistema: no controla ni sigue los exámenes \\
Docente & Ve solo sus grupos. Carga la lista de inscritos, registra el examen, habilita, elige aulas, asigna auxiliares y entrega los códigos QR. También controla el ingreso de sus exámenes \\
Auxiliar & Controla el ingreso en el examen y el aula donde el docente lo asignó. Puede quedar a cargo \propuesto \\
Estudiante & Sin cuenta. Recibe su código QR y lo presenta en la puerta \\
\hline
\end{tabular}
\end{center}

\section{Decisiones tomadas}

\begin{enumerate}
  \item La oferta académica se importa de una copia de los datos de GENDA. El
        \emph{scraper} se queda en GENDA.
  \item El sistema es un apoyo a la WebSIS, de uso institucional, no una herramienta
        opcional de cada docente.
  \item El período académico es una entidad y se detecta de lo que publica la
        universidad y de la fecha del sistema. Pueden estar vigentes dos a la vez.
  \item El grupo pertenece a un período. La asignatura, el docente, el aula y el
        estudiante no.
  \item El estudiante existe una sola vez, identificado por su código universitario. La
        inscripción lo une a un grupo.
  \item La lista de inscritos la carga el docente de su grupo o la administración. Nunca
        el estudiante.
  \item Si el estudiante ya existe, la carga lo reutiliza y avisa. Si sus datos no
        coinciden, queda como conflicto para la administración. El docente no edita los
        datos de un estudiante que ya existe.
  \item El examen es de la asignatura y abarca varios grupos, incluso de otros docentes.
  \item Las cuentas de docente se activan con el correo institucional y una contraseña
        temporal individual, de cambio obligatorio.
  \item El estudiante no tiene cuenta ni marca su propia presencia.
  \item El alcance son los exámenes de las materias, no el examen de ingreso.
  \item Solo entran las facultades cuya oferta se obtiene de forma automática: Ciencias
        y Tecnología, Economía, Humanidades y Arquitectura. Derecho y Odontología quedan
        fuera.
  \item El aula se usa como dato, sin capacidad: no hay fuente que la publique. El
        docente elige cualquiera de las aulas importadas.
  \item Aulas y docentes se ven del campus completo o por facultad, y cada facultad
        conserva su color.
  \item Un docente es una sola persona con una sola cuenta, aunque dicte en varias
        facultades.
  \item El docente ve todos sus exámenes del período en una lista, de todas sus
        asignaturas.
  \item Las pantallas muestran datos, estados y acciones, sin textos que expliquen el
        funcionamiento.
  \item No se usa la extensión de navegador para la WebSIS.
  \item Se quita el rol de responsable académico. No figura en el pliego, la consultora
        lo puso en duda y sus funciones quedaron cubiertas: el docente sigue su propio
        examen y el auxiliar, el aula que atiende. Los roles son tres: administrador,
        docente y auxiliar.
  \item El seguimiento en vivo es del docente y del auxiliar. El administrador se
        encarga de los datos del sistema, no de controlar ni de seguir los exámenes.
\end{enumerate}

\section{Consecuencias para el diseño, todavía sin confirmar}

De la sección 4 salen cambios que se propusieron y que el equipo aún no aprobó:

\begin{itemize}
  \item El docente opera el punto de control de sus propios exámenes. \propuesto
  \item El auxiliar no es un rol global: se asigna por examen y por aula, a cargo o de
        apoyo. \propuesto
  \item El auxiliar a cargo registra ingresos, expulsiones y situaciones excepcionales,
        pero no cambia la habilitación. \propuesto
  \item El docente tiene una vista de su examen con todas sus aulas, para cuando pasa de
        una a otra. \propuesto
  \item Un aula puede tener varios exámenes a la misma hora. En vez de impedirlo, el
        sistema avisa. \propuesto
  \item La puerta trabaja por aula y hora, no por examen: al identificar al estudiante,
        el sistema lo busca en todos los exámenes que hay en esa aula en ese momento.
        \propuesto
  \item Los reportes siguen siendo por examen, aunque los estudiantes rindan juntos.
        \propuesto
  \item Los auxiliares titulares se importan de los horarios, como los docentes.
        \propuesto
\end{itemize}

Dos criterios del \emph{backlog} dejan de aplicar por la decisión 13: en la HU-10, la
comparación de la capacidad asignada contra los habilitados; en la HU-14, que ningún
ambiente quede por encima de su capacidad. La HU-10 pide además que un ambiente no se
asigne a dos exámenes que se solapen, y eso cambia si se aprueba el aula compartida.

\section{Preguntas abiertas}

\begin{enumerate}
  \item ¿Cómo entra un auxiliar al sistema? Con cuenta propia, con un código de acceso
        por aula, o con las dos formas.
  \item Si dos docentes distintos comparten un aula, ¿alcanza con avisar o el que la
        tenía debe aceptar?
  \item En un examen conjunto de una materia, ¿quién puede cambiar la fecha, la hora y
        las aulas?
  \item ¿Qué información de estudiantes tiene realmente la administración?
  \item ¿Qué formato tiene la lista de inscritos que descarga el docente?
  \item ¿Se ubican en el mapa los auditorios de la FCyT?
  \item ¿Qué mapa base se usa debajo de los edificios?
\end{enumerate}

\section{Lo que se descartó}

\begin{center}
\begin{tabular}{@{}L{5.4cm}L{10.4cm}@{}}
\hline
\textbf{Opción} & \textbf{Por qué no} \\
\hline
Autoinscripción con código de curso & Deja entrar a quien no corresponde \\
El estudiante marca su presencia & Cualquiera podría marcarla por otro; el pliego pone la verificación en la puerta \\
Sistema opcional por docente & Sin lista confiable de inscritos no se cumple el requerimiento 9 \\
Contraseña inicial común a todos los docentes & Quien la conozca entra a cualquier cuenta antes del cambio \\
Extensión de navegador para la WebSIS & Trabaja con la cuenta y la contraseña de cada persona y no exporta listas \\
Capacidad de las aulas & No hay fuente que la publique \\
Derecho y Odontología & Derecho no tiene fuente web; los datos de Odontología son de un período anterior \\
Examen de ingreso a la universidad & No tiene materia, grupo ni código universitario \\
Rol de responsable académico & No figura en el pliego y lo que hacía lo cubren el docente y el auxiliar \\
\hline
\end{tabular}
\end{center}

\end{document}
